\documentclass[11pt]{article}

\usepackage[T1]{fontenc}
\usepackage[utf8]{inputenc}
\usepackage[spanish]{babel}
\usepackage{graphicx}
\usepackage{enumitem}

\parindent = 0cm
\setlist{itemsep=0.6em}

\usepackage{amsmath}
\usepackage{amssymb,amsfonts,latexsym,cancel}
\providecommand{\abs}[1]{\lvert#1\rvert}
\providecommand{\norm}[1]{\LVert#1\RVert}
\usepackage{multirow}
\usepackage[table,xcdraw]{xcolor}
\usepackage{booktabs}

% Márgenes limpios y compactos para todo el documento
\usepackage[lmargin=2cm,rmargin=2cm,top=2cm,bottom=2.5cm]{geometry}
\usepackage{xspace}

% ======================================================================
%  DEFINICIÓN DEL TÍTULO DEL TEMA
% ======================================================================
\newcommand{\titulotema}{Repaso Certamen 1}

% ======================================================================
%  CONFIGURACIÓN DE ENCABEZADOS Y PIES DE PÁGINA (fancyhdr)
% ======================================================================
\usepackage{fancyhdr}

% --- Estilo general: desde la página 2 en adelante ---
\pagestyle{fancy}
\fancyhf{}
\fancyhead[L]{Macroeconomía (580323)}
\fancyhead[R]{\titulotema}
\renewcommand{\headrulewidth}{0.5pt}
\fancyfoot[C]{\thepage}
\fancyfoot[R]{Semestre 2026-2}
\fancyfoot[L]{CMP/IGU/AFR}
\renewcommand{\footrulewidth}{0.5pt}

% --- Estilo exclusivo de la Página 1 ---
\fancypagestyle{primerapagina}{%
  \fancyhf{}
  \fancyhead[L]{\small FACULTAD DE INGENIERÍA\\DEPTO DE INGENIERÍA INDUSTRIAL \\ MACROECONOMÍA (580323)}
  \fancyhead[R]{\includegraphics[scale=0.13]{escudoudec.jpg}}
  \renewcommand{\headrulewidth}{0.9pt}
  \fancyfoot[C]{\thepage}
  \fancyfoot[R]{Semestre 2026-2}
  \fancyfoot[L]{CMP/IGU/AFR}
  \renewcommand{\footrulewidth}{0.5pt}
}

% Enlaces
\usepackage{url}
\usepackage{hyperref}
\hypersetup{
    colorlinks=true,
    linkcolor=blue,
    filecolor=magenta,
    urlcolor=cyan,
}

% ======================================================================
%  COMANDOS PERSONALIZADOS
% ======================================================================

\newcounter{ejercicio}
\newcommand{\ejercicio}[1][]{%
  \stepcounter{ejercicio}%
  \bigskip\noindent\textbf{Ejercicio \arabic{ejercicio}\ifx&#1&\relax.\else\ (#1).\fi}\ %
}

\newenvironment{partes}
  {\begin{enumerate}[label=(\alph*), leftmargin=1.5cm]}
  {\end{enumerate}}

\newenvironment{tabladatos}
  {\begin{center}}
  {\end{center}}

\newenvironment{recuadro}[1]{%
  \noindent\colorbox{black!8}{%
    \parbox{\dimexpr\linewidth-2\fboxsep\relax}{%
      \centering\Large\bfseries #1
    }%
  }\par\vspace{4mm}\noindent
}{\par}

\begin{document}
\thispagestyle{primerapagina}

\vspace*{0.05cm}

\begin{center}
  {\Large \textbf{Listado 4 -- Repaso -- Macroeconomía (580323)}}\\[0.15cm]
  {\large \titulotema}
\end{center}

\bigskip

\textbf{Profesor:} Cristian Mardones, Ph.D. \quad
\textbf{Ayudantes:} Ignacio Garrido U y Alex Fierro R.

\bigskip

%=======================================================================
% EJERCICIO -- Política Fiscal, Rigideces y Expectativas Racionales
%=======================================================================
\ejercicio[Certamen 2025]
Suponga que la econom\'ia se encuentra inicialmente en equilibrio. Ahora asuma que se est\'a discutiendo un incremento sustancial en la ley de presupuesto fiscal para el pr\'oximo a\~no. Analice y compare los efectos en el nivel de producci\'on y nivel de precios bajo los siguientes dos supuestos alternativos:

\begin{itemize}
    \item La curva de oferta agregada de corto plazo gira sobre su eje con el tiempo debido a rigideces nominales y ajustes lentos de precios.
\end{itemize}

Explique y grafique c\'omo se ajusta la econom\'ia desde el corto hasta el largo plazo y cu\'al es la efectividad del est\'imulo fiscal en cada escenario.

%=======================================================================
% EJERCICIO -- Ajuste Dinámico de Precios ante Caída de Importaciones
%=======================================================================
\ejercicio[Certamen 2023]
Suponga que la curva de oferta agregada chilena se puede caracterizar por la siguiente ecuaci\'on:
\[P_{t+1} = P_t \left[ 1 + \lambda (Y - Y^*) \right]\]
Explique y grafique el proceso de ajuste a trav\'es del tiempo con un modelo de oferta y demanda agregada si existe una importante ca\'ida en las importaciones. Preoc\'upese de etiquetar adecuadamente el gr\'afico, cualquier omisi\'on tendr\'a descuento.

%=======================================================================
% EJERCICIO -- Cuentas Nacionales, PIB Nominal y Cambio de Base
%=======================================================================
\ejercicio[Certamen 2021]
La producci\'on de bienes de una econom\'ia y sus respectivos precios desde 2016 hasta 2020 est\'an representados en la siguiente tabla. Los bienes de los sectores A, B y C son vendidos a consumidores finales, gobierno o al extranjero, mientras que el bien del sector D es utilizado como insumo para fabricar los otros tres bienes.

\begin{center}
\begin{tabular}{c c c c c c c c c}
\toprule
\textbf{Año} & $P_A$ & $Q_A$ & $P_B$ & $Q_B$ & $P_C$ & $Q_C$ & $P_D$ & $Q_D$ \\
\midrule
2016 & 8  & 50 & 12 & 10 & 16 & 80 & 6  & 10 \\
2017 & 13 & 55 & 16 & 11 & 19 & 84 & 9  & 12 \\
2018 & 14 & 58 & 18 & 14 & 22 & 82 & 12 & 13 \\
2019 & 13 & 53 & 15 & 22 & 25 & 83 & 15 & 18 \\
2020 & 30 & 54 & 16 & 22 & 28 & 88 & 18 & 19 \\
\bottomrule
\end{tabular}
\end{center}

\begin{partes}
    \item Determine la tasa de crecimiento del PIB nominal para cada a\~no y se\~nale si este indicador es \'util para reflejar el crecimiento de la econom\'ia.
    \item Demuestre c\'omo el cambio del a\~no base desde 2016 a 2020 modifica la tasa de crecimiento del PIB real (utilice porcentajes con un decimal, XX,X\%), luego mencione si existe una alternativa metodol\'ogica para solucionar el problema asociado al cambio del a\~no base.
\end{partes}

%=======================================================================
% EJERCICIO -- Auditoría de Cuentas Nacionales
%=======================================================================
\newpage
\ejercicio
Un analista junior le entrega el siguiente borrador de c\'alculo del PIB de un pa\'is para 2026 y le pide revisarlo antes de publicarlo. Su borrador \textbf{sum\'o} los siguientes conceptos:

\begin{enumerate}[label=\arabic*.]
    \item El valor de todas las ventas de trigo de agricultores a molinos, m\'as el valor de toda la harina vendida por los molinos a panader\'ias, m\'as el valor de todo el pan vendido a los hogares.
    \item El valor de todas las viviendas \textbf{usadas} compradas y vendidas durante el a\~no.
    \item El pago de intereses de la deuda p\'ublica del gobierno a los tenedores de bonos.
\end{enumerate}
Adem\'as, el analista \textbf{excluy\'o} deliberadamente del c\'alculo:
\begin{enumerate}[label=\arabic*., resume]
    \item La producci\'on de una f\'abrica de zapatos que no logr\'o vender un $15\%$ de su producci\'on, la cual qued\'o como inventario en bodega, argumentando que ``no se vendi\'o, as\'i que no genera actividad econ\'omica''.
\end{enumerate}
Y s\'i incluy\'o, sin dudarlo:
\begin{enumerate}[label=\arabic*., resume]
    \item La compra de un cami\'on nuevo por parte de una empresa de transporte de carga.
\end{enumerate}

\begin{partes}
    \item Para cada uno de los cinco \'items, indique si est\'a correctamente tratado, incorrectamente incluido, o incorrectamente excluido del PIB. Justifique cada caso con el concepto relevante (doble conteo, bienes usados, transferencias, inventarios como inversi\'on, etc.).
    \item Reformule correctamente el c\'alculo del \'item 1 (trigo $\to$ harina $\to$ pan), indicando qu\'e m\'etodo (gasto final o valor agregado) usar\'ia para evitar el error.
    \item Si el gobierno decide reclasificar el pago de intereses de la deuda p\'ublica como parte del gasto de gobierno ($G$), ?`qu\'e le dir\'ia usted como asesor t\'ecnico? Explique la diferencia conceptual entre un gasto y una transferencia.
\end{partes}

\end{document}